\section{Passaggio a 5G}

Nel lavoro di riferimento, la componente cellulare del sistema multi-tecnologia è basata su LTE C-V2X in una configurazione sotto copertura e assistita dalla rete \cite{segata2023multi-technology,3gpp-21914-1400}. Nel presente lavoro l'obiettivo non era cambiare il ruolo funzionale di questa terza interfaccia, ma aggiornarne la modellazione: il canale cellulare doveva continuare a comportarsi come una tecnologia gestita dall'infrastruttura e complementare a 802.11p e VLC, pur essendo realizzato con uno stack 5G piu` recente.

Simu5G rappresenta l'evoluzione naturale di SimuLTE nell'ecosistema OMNeT++ e rende possibile il passaggio a una modellazione 5G senza spezzare la continuita` con l'ambiente di simulazione gia` adottato \cite{virdis2014simulte,nardini2020simu5g}. Questo aspetto e` importante perche' la migrazione non consiste in una semplice sostituzione nominale di LTE con 5G: cambiano il framework di riferimento, l'organizzazione dei moduli di rete e parte dei parametri con cui vengono descritti accesso radio, mobilita` e transitori di connettivita`.

Dal punto di vista progettuale, il lavoro e` consistito nel trasferire dentro lo scenario multi-tecnologia una componente 5G gia` disponibile nei sotto-progetti di PLEXE, usandola come base invece di costruire da zero una nuova interfaccia cellulare. In questo modo e` stato possibile mantenere una logica coerente con il resto dell'ambiente: i beacon cooperativi continuano a essere trattati come traffico di platooning e la rete cellulare continua a rappresentare il terzo canale disponibile accanto a 802.11p e VLC.

Il punto chiave e` che il passaggio a 5G non modifica la logica superiore di \emph{SafeSwitch}. I follower continuano a valutare l'affidabilita` della componente cellulare insieme alle altre interfacce, e il suo contributo resta osservato attraverso metriche come il PDR, esattamente come avviene per gli altri canali. Quello che cambia e` il livello sottostante con cui questa interfaccia viene modellata: gli effetti di copertura, mobilita` tra celle e transitori radio non dipendono piu` dalla precedente pila LTE, ma dalla nuova modellazione fornita da Simu5G.
